\chapter{Muscle Models}
\label{ch:muscle-models}

\epigraph{%
  A muscle is not a motor. It is a spring with a memory,
  a damper with an agenda, and a force generator with a personality.
}{}

\section{The Hill-Type Muscle Model}

The Hill-type muscle model, originating from A.V.\ Hill's Nobel Prize-winning work on
muscle thermodynamics \cite{Hill1938}, is the standard lumped-parameter model used in
biomechanical simulation. It represents the musculotendon unit as a combination of
contractile, elastic, and damping elements.

\subsection{Model Architecture}

The standard Hill-type musculotendon model consists of three elements:

\begin{enumerate}
  \item \textbf{Contractile Element (CE)}: represents the active force-generating
        capacity of the muscle fibers. Its force depends on activation $a(t)$, fiber
        length $\ell^M$, and fiber velocity $v^M = \dot{\ell}^M$.
  \item \textbf{Parallel Elastic Element (PE)}: represents the passive elastic
        properties of the muscle belly (titin filaments, connective tissue). Produces
        force only when the muscle is stretched beyond its resting length.
  \item \textbf{Series Elastic Element (SE)}: represents the tendon, which is in
        series with the contractile element. Modeled as a nonlinear spring with a
        slack length $\ell_T^s$.
\end{enumerate}

\begin{center}
\begin{tikzpicture}[scale=1.0, >=Stealth]
  % Tendon (series elastic)
  \draw[thick] (0, 0) -- (1, 0);
  \draw[thick, decorate, decoration={zigzag, segment length=6, amplitude=3}]
    (1, 0) -- (3, 0) node[midway, above=8pt] {\small SE (tendon)};
  \draw[thick] (3, 0) -- (4, 0);

  % Junction point
  \filldraw (4, 0) circle (2pt);

  % Contractile element (top branch)
  \draw[thick] (4, 0) -- (4, 0.8) -- (5, 0.8);
  \draw[thick, fill=blue!10] (5, 0.4) rectangle (7, 1.2);
  \node at (6, 0.8) {\small CE};
  \draw[thick] (7, 0.8) -- (8, 0.8) -- (8, 0);

  % Parallel elastic element (bottom branch)
  \draw[thick] (4, 0) -- (4, -0.8) -- (5, -0.8);
  \draw[thick, decorate, decoration={zigzag, segment length=6, amplitude=3}]
    (5, -0.8) -- (7, -0.8) node[midway, below=8pt] {\small PE};
  \draw[thick] (7, -0.8) -- (8, -0.8) -- (8, 0);

  % Attachment points
  \draw[thick] (8, 0) -- (9, 0);
  \filldraw (0, 0) circle (2pt);
  \filldraw (9, 0) circle (2pt);

  % Labels
  \node[below] at (0, -0.1) {\small Origin};
  \node[below] at (9, -0.1) {\small Insertion};

  % Pennation angle
  \draw[thick, dashed, gray] (5, 0) -- (5, 0.4);
  \draw[gray, ->] (5, 0) arc (270:300:0.4);
  \node[gray, right] at (5.3, -0.2) {\small $\alpha$};
\end{tikzpicture}
\end{center}

\subsection{The Force-Length Relationship}

The isometric (zero velocity) force of the contractile element depends on the
normalized fiber length $\tilde{\ell} = \ell^M / \ell_0^M$:

\begin{definition}{Active Force-Length Curve}{fl-curve}
  The normalized active force-length relationship is:
  \begin{equation}
    f_L(\tilde{\ell}) = \exp\!\left(-\frac{(\tilde{\ell} - 1)^2}{\gamma}\right),
    \label{eq:ch3:force-length}
  \end{equation}
  where $\gamma \approx 0.45$ controls the width of the curve. Maximum force ($f_L = 1$)
  occurs at optimal fiber length ($\tilde{\ell} = 1$). Force decreases for both shorter
  and longer fiber lengths, reaching near-zero at approximately $\tilde{\ell} = 0.5$
  and $\tilde{\ell} = 1.5$.
\end{definition}

The physical basis for the force-length relationship is the sliding filament theory
\cite{HuxleySimmons1971}: at the optimal length, the maximum number of actin-myosin
cross-bridges can form. At shorter lengths, filament overlap reduces the number of
available binding sites. At longer lengths, the filaments pull apart.

\begin{lstlisting}[language=Python, caption={Hill-type muscle model implementation.}]
import numpy as np
from dataclasses import dataclass
from typing import Tuple

@dataclass
class MuscleParameters:
    """Parameters for a Hill-type musculotendon model."""
    F0: float          # Maximum isometric force (N)
    l0_M: float        # Optimal fiber length (m)
    v0: float          # Maximum contraction velocity (m/s)
    l_T_slack: float   # Tendon slack length (m)
    alpha0: float      # Pennation angle at optimal length (rad)
    gamma: float = 0.45  # Force-length curve width parameter
    Af: float = 0.25   # Force-velocity shape factor
    kPE: float = 4.0   # Passive elastic stiffness parameter
    e0: float = 0.6    # Passive elastic strain at F0

    @property
    def v_max(self) -> float:
        """Maximum shortening velocity in fiber lengths/s."""
        return self.v0 / self.l0_M


def force_length_active(l_tilde: float, gamma: float = 0.45) -> float:
    """Active force-length relationship.

    Parameters
    ----------
    l_tilde : float
        Normalized fiber length (l_M / l0_M).
    gamma : float
        Width parameter of the Gaussian curve.

    Returns
    -------
    float
        Normalized active force (0 to 1).

    References
    ----------
    Thelen, D.G. (2003). Adjustment of muscle mechanics model
    parameters to simulate dynamic contractions in older adults.
    J. Biomechanical Engineering, 125(1), 70-77.
    """
    return np.exp(-((l_tilde - 1.0) ** 2) / gamma)


def force_length_passive(l_tilde: float, kPE: float = 4.0,
                         e0: float = 0.6) -> float:
    """Passive force-length relationship (exponential model).

    Parameters
    ----------
    l_tilde : float
        Normalized fiber length.
    kPE : float
        Stiffness parameter.
    e0 : float
        Strain at which passive force equals F0.

    Returns
    -------
    float
        Normalized passive force.
    """
    if l_tilde <= 1.0:
        return 0.0
    return (np.exp(kPE * (l_tilde - 1.0) / e0 - kPE) - 1.0) / \
           (np.exp(kPE) - 1.0)


def force_velocity(v_tilde: float, a: float, Af: float = 0.25,
                   Flen: float = 1.4) -> float:
    """Force-velocity relationship (Thelen 2003 model).

    Parameters
    ----------
    v_tilde : float
        Normalized fiber velocity (v_M / v_max).
        Negative = shortening (concentric).
    a : float
        Activation level (0 to 1).
    Af : float
        Shape factor for concentric portion.
    Flen : float
        Maximum normalized eccentric force.

    Returns
    -------
    float
        Normalized force multiplier.
    """
    if v_tilde <= 0:
        # Concentric (shortening)
        return (1.0 + v_tilde) / (1.0 - v_tilde / (Af * a + 0.01))
    else:
        # Eccentric (lengthening)
        return (1.0 + v_tilde * Flen / (0.04 + a))  / \
               (1.0 + v_tilde / (0.04 + a))


def muscle_force(
    activation: float,
    l_M: float,
    v_M: float,
    params: MuscleParameters
) -> Tuple[float, float, float]:
    """Compute total musculotendon force.

    Parameters
    ----------
    activation : float
        Neural activation level (0 to 1).
    l_M : float
        Current muscle fiber length (m).
    v_M : float
        Current muscle fiber velocity (m/s, negative = shortening).
    params : MuscleParameters
        Muscle parameters.

    Returns
    -------
    Tuple[float, float, float]
        (total_force, active_force, passive_force) in Newtons.
    """
    # Normalize
    l_tilde = l_M / params.l0_M
    v_tilde = v_M / params.v0

    # Force components
    f_active = (activation
                * force_length_active(l_tilde, params.gamma)
                * force_velocity(v_tilde, activation, params.Af))
    f_passive = force_length_passive(l_tilde, params.kPE, params.e0)

    # Pennation angle (changes with fiber length)
    cos_alpha = np.cos(params.alpha0) if l_tilde > 0.1 else 1.0

    # Total force along tendon
    F_active = f_active * params.F0 * cos_alpha
    F_passive = f_passive * params.F0 * cos_alpha
    F_total = F_active + F_passive

    return F_total, F_active, F_passive


# Example: Biceps brachii
biceps = MuscleParameters(
    F0=624.3,        # N (Arnold et al. 2010)
    l0_M=0.1157,     # m
    v0=10 * 0.1157,  # 10 fiber lengths/s
    l_T_slack=0.2723, # m
    alpha0=0.0        # parallel-fibered
)

# Compute force at optimal length, half activation, isometric
F_total, F_active, F_passive = muscle_force(
    activation=0.5,
    l_M=biceps.l0_M,  # optimal length
    v_M=0.0,           # isometric
    params=biceps
)
print(f"Biceps at 50% activation, optimal length, isometric:")
print(f"  Active force:  {F_active:.1f} N")
print(f"  Passive force: {F_passive:.1f} N")
print(f"  Total force:   {F_total:.1f} N")
\end{lstlisting}

\subsection{The Force-Velocity Relationship}

The second critical property of the contractile element is the force-velocity
relationship, first characterized by Hill (1938) \cite{Hill1938}.

\begin{principle}{The Force-Velocity Asymmetry}{fv-asymmetry}
  \begin{itemize}
    \item \textbf{Concentric (shortening) contraction}: force \emph{decreases} with
          increasing shortening velocity. At maximum shortening velocity $v_0$, the
          muscle produces zero force. The relationship is approximately hyperbolic:
          $(F + a)(v + b) = (F_0 + a) b$.
    \item \textbf{Eccentric (lengthening) contraction}: force \emph{increases} beyond
          the maximum isometric force, typically reaching $1.4$--$1.8 \times F_0$.
          This is why muscles are ``stronger'' when resisting a load than when lifting it.
    \item \textbf{Isometric ($v = 0$)}: force equals the product of activation and
          the force-length relationship.
  \end{itemize}
  The asymmetry between concentric and eccentric force production is unique to
  biological actuators. No common engineering actuator exhibits this behavior.
\end{principle}

\subsection{Activation Dynamics}

The neural command from the central nervous system does not instantly produce
muscle force. The activation dynamics models the excitation-contraction coupling:
\begin{equation}
  \dot{a}(t) = \frac{u(t) - a(t)}{\tau(u, a)},
  \label{eq:ch3:activation-dynamics}
\end{equation}
where $u(t) \in [0, 1]$ is the neural excitation, $a(t) \in [0, 1]$ is the
activation level, and $\tau$ is a time constant that differs for activation
and deactivation:
\begin{equation}
  \tau(u, a) = \begin{cases}
    \tau_{\text{act}}(0.5 + 1.5a) & \text{if } u > a \text{ (activation)}, \\
    \tau_{\text{deact}} / (0.5 + 1.5a) & \text{if } u \leq a \text{ (deactivation)}.
  \end{cases}
  \label{eq:ch3:tau-activation}
\end{equation}
Typical values: $\tau_{\text{act}} = 10$--$15$~ms, $\tau_{\text{deact}} = 40$--$60$~ms
\cite{Winters1995}. The asymmetry between activation and deactivation time constants
has important consequences for motor control: it is much easier to ``turn on'' a
muscle quickly than to ``turn it off.''

\section{Tendon Mechanics}

The tendon connects the muscle belly to bone. It is modeled as a nonlinear elastic
element with a resting (slack) length $\ell_T^s$:
\begin{equation}
  F_T(\varepsilon_T) = \begin{cases}
    0 & \varepsilon_T \leq 0, \\
    F_0 \cdot \frac{e^{k_T \varepsilon_T} - 1}{e^{k_T \varepsilon_{0T}} - 1} &
    \varepsilon_T > 0,
  \end{cases}
  \label{eq:ch3:tendon-force}
\end{equation}
where $\varepsilon_T = (\ell_T - \ell_T^s) / \ell_T^s$ is the tendon strain and
$\varepsilon_{0T} \approx 0.033$ is the strain at maximum isometric force. The
tendon stiffness parameter $k_T$ is typically chosen so that the tendon strain is
approximately 3.3\% at $F_0$ \cite{Zajac1989}.

\begin{clinicalbox}{Energy Storage in Tendons}{tendon-energy}
  Tendons are not merely passive connectors---they are \emph{energy storage devices}.
  During running, the Achilles tendon stores approximately 35~J of elastic energy
  per stride, which is returned during push-off. This reduces the metabolic cost
  of running by approximately 50\% compared to a hypothetical system with rigid
  tendons \cite{Alexander1991}.

  For the golf swing, the tendons of the wrist and forearm muscles store energy
  during the early downswing (when the wrist is cocked) and release it during
  the uncocking phase, contributing to clubhead speed.
\end{clinicalbox}

\section{Musculotendon Equilibrium}

The muscle fiber and tendon must satisfy force equilibrium at all times. Assuming
the muscle fiber acts at a pennation angle $\alpha$ to the tendon line of action:
\begin{equation}
  F_T = (F_{\text{CE}} + F_{\text{PE}}) \cos\alpha,
  \label{eq:ch3:mt-equilibrium}
\end{equation}
where:
\begin{align}
  F_{\text{CE}} &= a \cdot f_L(\tilde{\ell}^M) \cdot f_V(\tilde{v}^M) \cdot F_0, \\
  F_{\text{PE}} &= f_{\text{PE}}(\tilde{\ell}^M) \cdot F_0.
\end{align}

The musculotendon length $\ell^{MT}$ is the sum of tendon length and the
projection of fiber length:
\begin{equation}
  \ell^{MT} = \ell_T + \ell^M \cos\alpha.
  \label{eq:ch3:mt-length}
\end{equation}

Given a known musculotendon length (from the skeletal geometry) and activation,
the equilibrium condition \eqref{eq:ch3:mt-equilibrium} and the constraint
\eqref{eq:ch3:mt-length} determine the fiber length and velocity.

\section{The Complete Hill Model as a Dynamical System}

Combining all elements, the Hill muscle model can be expressed as a dynamical
system with two state variables per muscle:
\begin{align}
  \dot{a} &= \frac{u - a}{\tau(u, a)}, \label{eq:ch3:state-activation} \\
  \dot{\ell}^M &= v^M(a, \ell^M, \ell^{MT}), \label{eq:ch3:state-fiber}
\end{align}
where $v^M$ is obtained by inverting the force-velocity relationship at the
equilibrium fiber force. The input is the neural excitation $u(t) \in [0, 1]$
and the ``disturbance'' is the musculotendon path length $\ell^{MT}(\bm{q})$,
which depends on the joint configuration.

\begin{principle}{Muscle as a Nonlinear Control System}{muscle-control}
  Each muscle is itself a nonlinear dynamical system with:
  \begin{itemize}
    \item \textbf{State}: $(a, \ell^M) \in [0,1] \times \Reals_{>0}$
    \item \textbf{Input}: neural excitation $u \in [0,1]$
    \item \textbf{Output}: tendon force $F_T$
    \item \textbf{Disturbance}: musculotendon path length $\ell^{MT}(\bm{q})$
  \end{itemize}
  The muscle transforms a scalar neural command into a scalar force, but the
  transformation is highly nonlinear, state-dependent, and history-dependent.
  The entire musculoskeletal system is a cascade: neural commands $\to$ muscle
  dynamics $\to$ tendon forces $\to$ joint torques $\to$ skeletal dynamics.
\end{principle}

\section*{Exercises}

\begin{enumerate}
  \item \textbf{Force-Length Curve.}
        Plot the active force-length curve for $\gamma \in \{0.3, 0.45, 0.6\}$.
        For each, determine the ``useful range'' (the interval of fiber lengths
        where $f_L > 0.5$). How does the width parameter affect the operating range?

  \item \textbf{Force-Velocity Curve.}
        Plot the complete force-velocity curve (both concentric and eccentric) for
        the biceps parameters given in the code listing. Mark the maximum eccentric
        force and the maximum shortening velocity.

  \item \textbf{Activation Dynamics.}
        Simulate the activation dynamics for a square-wave neural excitation input
        ($u = 1$ for 200~ms, then $u = 0$). Plot the activation response with
        $\tau_{\text{act}} = 15$~ms and $\tau_{\text{deact}} = 50$~ms. How long does
        it take for activation to reach 95\% of maximum? How long for deactivation?

  \item \textbf{Tendon Compliance.}
        For the Achilles tendon (slack length $\ell_T^s = 0.25$~m, max force 4000~N),
        compute the tendon stretch at 50\%, 75\%, and 100\% of maximum force. Does
        the tendon behave approximately linearly over this range?

  \item \textbf{Isometric Force Surface.}
        Create a 3D surface plot of muscle force as a function of normalized fiber
        length ($\tilde{\ell}$) and activation ($a$) at zero velocity (isometric).
        Identify the point of maximum force production.

  \item \textbf{(Programming.)} Implement a complete musculotendon simulation: given
        a prescribed musculotendon length trajectory $\ell^{MT}(t)$ and neural
        excitation $u(t)$, integrate the state equations \eqref{eq:ch3:state-activation}
        and \eqref{eq:ch3:state-fiber} and plot the resulting tendon force $F_T(t)$.

  \item \textbf{(Programming.)} Implement a two-muscle antagonist pair (e.g., biceps
        and triceps) acting on a single-DOF elbow joint. Given co-contraction
        ($u_{\text{biceps}} = u_{\text{triceps}} = 0.5$), compute the net joint torque
        and the joint stiffness. Verify that co-contraction increases stiffness without
        changing the equilibrium position.

  \item \textbf{(Research.)} Compare the Thelen (2003) muscle model with the
        Millard et al.\ (2013) model used in OpenSim 4.x. What are the key differences
        in the force-velocity and passive force models? Which is more physiologically
        accurate?
\end{enumerate}
